\documentclass[11pt,a4paper]{article}
\usepackage{amsmath,amssymb}
\usepackage{graphicx}
\usepackage{booktabs}
\usepackage{hyperref}
\usepackage[margin=2.5cm]{geometry}
\usepackage{caption}
\usepackage{subcaption}
\usepackage{placeins}

\title{Stage 4A — Multi-Neighbour Model Selection Report}
\author{MSc Student — BIOE70025}
\date{\today}

% Path to figures directory (set per user)
\newcommand{\figdir}{../experiments/logs/stage4a_model_selection/20260616_113914_final_evaluation/figures}

\begin{document}
\maketitle

% ======================================================================
\section{Introduction}
% ======================================================================

This stage report documents the Stage 4A multi-neighbour simulation-based
model selection for synchronisation algorithms prior to hardware-in-the-loop
(HIL) testing.  Five candidate models were evaluated under identical
decentralised multi-agent conditions using a pure-software simulator.  The
goal is to select a primary and a secondary HIL candidate based on strict
zero-lag synchronisation performance, robustness across topologies and
frequency sets, N=5 scalability, hardware readiness, and interpretability.

% ======================================================================
\section{Candidate Models}
% ======================================================================

All models produce flash outputs and are evaluated with common flash-based
synchronisation metrics.  The PCO and EAPF variants are event-driven: each
agent updates its state only when neighbour flash events are detected.
Kuramoto uses continuous phase coupling in simulation and would require
phase reconstruction from flash timestamps for hardware implementation.
Locked parameters were selected via a 60\,s $\times$ 5-repeat parameter sweep
(seeds 1000--1004) across four frequency sets and three topologies.

% --- Table 1 (locked params) placed after Candidate Models intro ---
% ======================================================================
% Stage 4A Model Selection — Tables
% ======================================================================

% --- Table 1: Locked Model Parameters ---
\begin{table}[htbp]
\centering
\caption{Locked model parameters selected via 60\,s $\times$ 5-repeat
         parameter-locking sweep (seeds 1000--1004).}
\label{tab:locked}
\begin{tabular}{@{}lll@{}}
\toprule
\textbf{Model} & \textbf{Key Parameters} & \textbf{Variant} \\
\midrule
Kuramoto & $K = 5.0$ & baseline \\
PCO simple & additive phase, $\varepsilon = 0.10$, $\tau_{\text{ref}} = 0.05$\,s & simple \\
PCO adaptive PRC & biphasic sine PRC, $\varepsilon = 0.10$, phase delay on,
                  $\Delta\phi_{\max} = 0.20$ & adaptive\_prc \\
EAPF tracker & $g_p = 0.40$, $g_f = 0.15$ & tracker \\
EAPF consensus & $g_p = 0.02$, $g_f = 0.02$, $\alpha_{p,f} = 0.2$,
                 $\Delta\theta_{\max} = 0.2$\,rad, $\Delta f_{\max} = 0.05$\,Hz & consensus \\
\bottomrule
\end{tabular}
\end{table}

% --- Table 2: Final Weighted Scores ---
\begin{table}[htbp]
\centering
\caption{Final weighted model evaluation scores (N=3: 1200 trials, N=5: 300 trials,
         seeds 2000--2019 / 3000--3009).}
\label{tab:scores}
\begin{tabular}{@{}lrrrrrrr@{}}
\toprule
\textbf{Model} & \textbf{N3 Sync} & \textbf{N3 Conv} & \textbf{N3 Robust}
               & \textbf{N5 Scale} & \textbf{HW} & \textbf{Int.} & \textbf{Total} \\
\midrule
EAPF consensus & 1.000 & 0.984 & 1.000 & 0.933 & 0.800 & 0.800 & \textbf{0.938} \\
Kuramoto       & 0.902 & 0.891 & 0.767 & 0.427 & 0.850 & 0.900 & 0.801 \\
EAPF tracker   & 0.763 & 0.752 & 0.335 & 0.597 & 0.700 & 0.700 & 0.656 \\
PCO adaptive   & 0.350 & 0.467 & 0.100 & 0.000 & 0.650 & 0.750 & 0.363 \\
PCO simple     & 0.283 & 0.370 & 0.100 & 0.000 & 0.750 & 0.900 & 0.358 \\
\bottomrule
\end{tabular}
\end{table}

% --- Table 3: Repeatability Comparison ---
\begin{table}[htbp]
\centering
\caption{Repeatability across two independent final evaluations (seeds
         2000/3000 vs.\ 4000/5000).}
\label{tab:repeatability}
\begin{tabular}{@{}lrrr@{}}
\toprule
\textbf{Model} & \textbf{Run 1} & \textbf{Run 2} & \textbf{$\Delta$} \\
\midrule
EAPF consensus & 0.9375 & 0.9375 & 0.0000 \\
Kuramoto       & 0.8008 & 0.8020 & $+$0.0012 \\
EAPF tracker   & 0.6563 & 0.6602 & $+$0.0039 \\
PCO adaptive   & 0.3625 & 0.3625 & 0.0000 \\
PCO simple     & 0.3579 & 0.3579 & 0.0000 \\
\bottomrule
\end{tabular}
\end{table}

% --- Table 4: Final Decision Summary ---
\begin{table}[htbp]
\centering
\caption{Final Stage 4A model-selection decision summary.}
\label{tab:decision}
\begin{tabular}{@{}llp{6.5cm}@{}}
\toprule
\textbf{Rank} & \textbf{Model} & \textbf{Decision} \\
\midrule
1 & EAPF consensus & \textbf{Primary HIL candidate} \\
2 & Kuramoto & Secondary comparison model \\
3 & EAPF tracker & Retained — works for near-identical oscillators \\
4 & PCO adaptive PRC & Retained — reduced extra flashes vs.\ simple PCO \\
5 & PCO simple & Retained — simplest firefly-inspired baseline \\
\bottomrule
\end{tabular}
\end{table}


% ----------------------------------------------------------------------
\subsection{Kuramoto (continuous phase coupling)}
% ----------------------------------------------------------------------

\textbf{Working principle \cite{kuramoto1975self}.}
Each agent $i$ continuously observes the phase of its visible neighbours
and adjusts its own phase velocity according to the Kuramoto ODE:
\begin{equation}
  \dot\theta_i = \omega_i + \frac{K}{N_i}\sum_{j\in\mathcal{N}_i}
                 \sin(\theta_j - \theta_i)
\end{equation}
where $\mathcal{N}_i$ is the neighbour set of agent $i$, $N_i = |\mathcal{N}_i|$,
and $K$ is the coupling strength.

\textbf{Locked parameter:} $K = 5.0$.

\textbf{Strengths:} mathematically well-understood; achieves 91.7\% zero-lag
success in N=3; simple to implement on hardware using phase reconstruction
from flash timestamps.

\textbf{Limitations:} requires continuous phase observability which the
camera pipeline approximates via flash-to-phase estimation; N=5 scalability
limited (33.3\% zero-lag success).

% ----------------------------------------------------------------------
\subsection{PCO Simple (event-driven pulse coupling)}
% ----------------------------------------------------------------------

\textbf{Working principle \cite{peskin1975mathematical,mirollo1990synchronization}.}
Each agent maintains a normalised phase $\phi \in [0,1)$.  On detecting a
neighbour flash, the phase is additively advanced:
\begin{equation}
  \phi \leftarrow \min(1.0,\; \phi + \varepsilon)
\end{equation}
When $\phi \geq 1.0$ the agent fires and resets to 0, entering a refractory
period.  This is the simplest firefly-inspired pulse-coupled oscillator,
derived from the Mirollo--Strogatz model.

\textbf{Locked parameters:} additive phase, $\varepsilon = 0.10$,
$\tau_{\text{ref}} = 0.05$\,s.

\textbf{Strengths:} directly event-driven; matches the camera detection
pipeline; conceptually simplest firefly model.

\textbf{Limitations:} 25.0\% zero-lag success at N=3; 0\% at N=5;
flash count ratio 1.308 (extra flashes in heterogeneous conditions).

% ----------------------------------------------------------------------
\subsection{PCO Adaptive PRC (phase response curve variant)}
% ----------------------------------------------------------------------

\textbf{Working principle \cite{mirollo1990synchronization,tyrrell2006fireflies}.}
Extends PCO simple with a biphasic sine phase response curve:
\begin{equation}
  \Delta\phi = \varepsilon \cdot \sin\bigl(2\pi(\phi - \phi_{\text{ref}})\bigr)
\end{equation}
Phase delay is enabled (negative corrections allowed), and maximum correction
is rate-limited.  Optional frequency adaptation is available but was not used
in the locked variant.

\textbf{Locked parameters:} biphasic sine PRC, $\varepsilon = 0.10$,
phase delay enabled, $\Delta\phi_{\max} = 0.20$.

\textbf{Strengths:} lower flash count ratio (1.175 vs.\ 1.308 for simple PCO);
better 1:1 flash lock rate (75\% vs.\ 42\%).

\textbf{Limitations:} 25.0\% zero-lag success at N=3; 0\% at N=5;
still limited by excitatory pulse coupling under strong heterogeneity.

% ----------------------------------------------------------------------
\subsection{EAPF Tracker (fixed-leader phase/frequency tracker)}
% ----------------------------------------------------------------------

\textbf{Working principle \cite{gardner2005phaselock}.}
On each detected leader flash, the agent computes the signed phase error
and applies independent phase and frequency corrections:
\begin{align}
  \theta &\leftarrow \operatorname{wrap}_{2\pi}(\theta + g_p \cdot e_\theta) \\
  f      &\leftarrow \operatorname{clamp}(f + g_f \cdot e_\theta / (2\pi),
                                          f_{\min}, f_{\max})
\end{align}
where $e_\theta = \operatorname{wrap}_\pi(0 - \theta)$.

\textbf{Locked parameters:} $g_p = 0.40$, $g_f = 0.15$.

\textbf{Strengths:} simple engineering tracker; 76.2\% zero-lag success at N=3;
works well for near-identical oscillators.

\textbf{Limitations:} designed for fixed-leader tracking; less robust in
mutual multi-neighbour mode (flash count ratio 1.116).

% ----------------------------------------------------------------------
\subsection{EAPF Consensus (event-based phase/frequency consensus)}
% ----------------------------------------------------------------------

\textbf{Working principle \cite{gardner2005phaselock,olfati2007consensus}.}
Each agent maintains per-neighbour state estimates (period, frequency, phase)
updated \emph{only} from flash timestamps.  Consensus corrections use the mean
phase and frequency errors across all visible neighbours:
\begin{align}
  e_\theta &= \frac{1}{|\mathcal{N}_i|}\sum_{j\in\mathcal{N}_i}
             \operatorname{wrap}_\pi(\hat\theta_j - \theta_i) \\
  e_f      &= \frac{1}{|\mathcal{N}_i|}\sum_{j\in\mathcal{N}_i}
             (\hat f_j - f_i)
\end{align}
where $\hat\theta_j$ and $\hat f_j$ are the estimated neighbour phase and
frequency.  Corrections are low-pass filtered and rate-limited.

\textbf{Locked parameters:} $g_p = 0.02$, $g_f = 0.02$,
$\alpha_p = \alpha_f = 0.2$, $\Delta\theta_{\max} = 0.2$\,rad,
$\Delta f_{\max} = 0.05$\,Hz.

\textbf{Strengths:} 100\% zero-lag success at N=3 \emph{and} N=5; flash count
ratio 1.000; nearest-flash timing error 6.6\,ms; frequency spread 0.0\,Hz;
only model achieving perfect strict synchronisation across all conditions.

\textbf{Limitations:} very low gains require 60\,s for full convergence;
more stateful than simpler models (per-neighbour estimates).

% ======================================================================
\section{Evaluation}
% ======================================================================

% ----------------------------------------------------------------------
\subsection{Evaluation Setup}
% ----------------------------------------------------------------------

The multi-agent simulator runs $N$ agents at fixed time-step $dt = 0.01$\,s.
Each agent observes neighbour flash events according to a graph-based topology.

\paragraph{N=3 main evaluation.}
The primary HIL-aligned evaluation uses 3 agents with topologies
\texttt{all\_to\_all}, \texttt{chain}, and \texttt{directed\_chain}, and four
frequency sets: identical $[2.0, 2.0, 2.0]$\,Hz, near-identical
$[1.9, 2.0, 2.1]$\,Hz, moderate heterogeneity $[1.8, 2.0, 2.2]$\,Hz, and
strong heterogeneity $[1.5, 2.0, 2.3]$\,Hz.  For each model, this gives
$4 \times 3 \times 20 = 240$ trials.  Across all five models, the N=3
evaluation comprises $5 \times 4 \times 3 \times 20 = 1200$ trials, each
of 60\,s duration with independent random seeds (2000--2019).

\paragraph{N=5 scalability evaluation.}
A second evaluation with 5 agents uses three local topologies
(\texttt{local\_ring\_5}, \texttt{chain\_5}, \texttt{local\_degree\_2\_3})
and two frequency sets (near-identical and strong heterogeneity), 10 repeats
per combination (seeds 3000--3009), producing $5 \times 3 \times 2 \times 10
= 300$ trials total.

\paragraph{Repeatability.}
A second independent run with seeds 4000--4019 (N=3) and 5000--5009 (N=5)
was performed with all locked parameters unchanged.

% ----------------------------------------------------------------------
\subsection{Metrics}
% ----------------------------------------------------------------------

\paragraph{Synchronisation quality.}
Zero-lag group sync requires phase synchronisation (pairwise timing error
$< 0.10$\,s and order parameter $R > 0.85$), frequency lock (spread
$< 0.05$\,Hz), and 1:1 flash lock (flash count ratio $\leq 1.2$).
Phase-locked group sync relaxes the in-phase constraint and allows stable
non-zero offsets (offset jitter $< 0.05$\,s).

\paragraph{Convergence and steady-state.}
Time-to-sync is the first time the 5-cycle sustained criterion is met.
Steady-state metrics (final 10\,s) include mean pairwise timing error,
offset jitter, frequency spread, and order parameter.

\paragraph{Flash behaviour.}
Flash count ratio is the maximum final-window flash count divided by the
minimum, measuring extra-flash or harmonic-locking behaviour.

% ----------------------------------------------------------------------
\subsection{Weighted Evaluation Matrix}
% ----------------------------------------------------------------------

The final model score is a weighted sum of six categories:
N=3 synchronisation performance (0.30), N=3 convergence/quality (0.15),
N=3 robustness across conditions (0.15), N=5 scalability (0.15),
hardware readiness (0.15, expert-assigned), and interpretability (0.10,
expert-assigned).

\FloatBarrier

% ======================================================================
\section{Results and Discussion}
% ======================================================================

% ----------------------------------------------------------------------
\subsection{N=3 Performance}
% ----------------------------------------------------------------------

EAPF consensus achieved 100\% zero-lag group sync across all 240 N=3 trials
for that model (4 frequency sets $\times$ 3 topologies $\times$ 20 repeats).
Kuramoto achieved 91.7\% zero-lag success, with failures concentrated in
chain and directed\_chain topologies under strong heterogeneity.  EAPF
tracker achieved 76.2\%, and both PCO variants achieved 25.0\%.

The flash count ratio (Table~\ref{tab:scores}) clearly separates models:
EAPF consensus and Kuramoto maintain FCR $= 1.0$; EAPF tracker (1.12) and
PCO adaptive PRC (1.18) show moderate extra flashes; PCO simple (1.31) shows
significant extra flashes indicating repeated near-simultaneous group flash
clusters with additional isolated flashes from faster agents.

% ----------------------------------------------------------------------
\subsection{N=5 Scalability}
% ----------------------------------------------------------------------

N=5 sharply separates the models.  EAPF consensus is the only model
maintaining 100\% zero-lag sync.  EAPF tracker drops to 58.3\%, Kuramoto to
33.3\%, and both PCO variants to 0\%.  The N=5 results confirm that consensus
PLL with low gains scales to larger agent counts under local-neighbourhood
topologies, while the other models degrade.

% ----------------------------------------------------------------------
\subsection{Model Ranking}
% ----------------------------------------------------------------------

Table~\ref{tab:scores} presents the final weighted scores.  EAPF consensus
ranks first (0.938), followed by Kuramoto (0.801), EAPF tracker (0.656--0.660),
PCO adaptive PRC (0.363), and PCO simple (0.358).

% ----------------------------------------------------------------------
\subsection{Repeatability}
% ----------------------------------------------------------------------

The second independent evaluation (seeds 4000/5000) produced nearly identical
scores (Table~\ref{tab:repeatability}): the maximum delta was 0.0039 (EAPF
tracker, $+$0.0039).  All models are within the ``very stable'' band
($|\Delta| < 0.03$).  The ranking is unchanged: EAPF consensus remains first,
Kuramoto second.

% ----------------------------------------------------------------------
\subsection{Interpretation}
% ----------------------------------------------------------------------

\paragraph{Why EAPF consensus outperforms Kuramoto.}
EAPF consensus is a fully event-driven distributed frequency/phase agreement
protocol \cite{gardner2005phaselock,olfati2007consensus}.  It converges slowly
(low gains) but robustly, without overshoot or correction-induced extra
flashes.  At 60\,s, even the 1.5--2.3\,Hz 0.8\,Hz frequency spread is
resolved to 0.0\,Hz final spread.  Kuramoto's continuous coupling is stronger
but less robust under directed/sparse topologies.

\paragraph{Why Kuramoto remains strong but less scalable.}
Kuramoto \cite{kuramoto1975self} achieves 91.7\% zero-lag at N=3 but drops to
33.3\% at N=5 because continuous coupling weakens as the number of indirect
neighbours increases and the effective coupling per link becomes diluted.

\paragraph{Why PCO variants remain limited.}
Excitatory phase-advance coupling \cite{mirollo1990synchronization} cannot
decelerate an oscillator that is already too fast.  Under strong heterogeneity,
the 2.3\,Hz agent consistently emits extra flashes.  The adaptive PRC
(biphasic sine with phase delay) reduces but does not eliminate this
behaviour.  This is consistent with the firefly-inspired synchronisation
limitations discussed in \cite{tyrrell2006fireflies}.

% ----------------------------------------------------------------------
\subsection{Flash Raster Interpretation}
% ----------------------------------------------------------------------

Figure~\ref{fig:raster_full} shows representative flash rasters over the full
0--60\,s interval for the all\_to\_all, strong heterogeneity condition.
Figure~\ref{fig:raster_zoom} provides a 0--10\,s zoom to reveal early-time
convergence dynamics.

EAPF consensus (panel b) shows agents converging to a single synchronised
flash cluster by approximately $t = 20$\,s, with perfect 1:1 locking
thereafter.  Kuramoto (panel a) shows rapid convergence within the first
3--5\,s.  EAPF tracker (panel c) shows slower convergence with some residual
offset.  PCO adaptive PRC (panel d) and PCO simple (panel e) show clustered
but non-1:1 flash patterns, with the fastest agent consistently emitting
additional flashes between group events.

Raster plots should be interpreted together with quantitative metrics:
the 0--60\,s raster shows overall structure; early-time offsets may be
visually compressed; the 0--10\,s zoom reveals initial convergence;
time-to-sync and timing-error metrics provide quantitative confirmation.

% ======================================================================
% Stage 4A Model Selection — Figures
% Figure paths are relative to the figures/ directory inside the final
% evaluation output folder.  Set \figdir in the main document.
% ======================================================================

% --- Figure 1: N=3 Zero-Lag Success Heatmap ---
\begin{figure}[htbp]
\centering
\includegraphics[width=0.95\textwidth]{\figdir/fig1_n3_zero_lag_success_heatmap.png}
\caption{N=3 zero-lag group sync success rate by model and frequency set
         (all topologies pooled).}
\label{fig:n3_heatmap}
\end{figure}

% --- Figure 7: N=5 Scalability Heatmap ---
\begin{figure}[htbp]
\centering
\includegraphics[width=0.85\textwidth]{\figdir/fig7_n5_scalability_success_heatmap.png}
\caption{N=5 scalability — zero-lag success rate by model and frequency set
         (all topologies pooled).}
\label{fig:n5_heatmap}
\end{figure}

% --- Figure 4: Flash Count Ratio ---
\begin{figure}[htbp]
\centering
\includegraphics[width=0.85\textwidth]{\figdir/fig4_n3_flash_count_ratio_comparison.png}
\caption{N=3 mean flash count ratio by model.  The dashed line at 1.2 marks
         the 1:1 flash lock threshold.}
\label{fig:fcr}
\end{figure}

% --- Figure 8: Weighted Scores ---
\begin{figure}[htbp]
\centering
\includegraphics[width=0.85\textwidth]{\figdir/fig8_weighted_model_scores.png}
\caption{Total weighted model evaluation scores.}
\label{fig:total}
\end{figure}

% --- Figure 9: Category Breakdown ---
\begin{figure}[htbp]
\centering
\includegraphics[width=0.95\textwidth]{\figdir/fig9_category_score_breakdown.png}
\caption{Category score breakdown: N3 Sync (0.30), N3 Conv (0.15),
         N3 Robust (0.15), N5 Scale (0.15), HW Ready (0.15), Interp (0.10).}
\label{fig:categories}
\end{figure}

% --- Figure 10: Full 0-60 s Flash Rasters ---
\begin{figure}[htbp]
\centering
\includegraphics[width=0.95\textwidth]{\figdir/fig10_representative_flash_rasters_full_0_60s.png}
\caption{Representative flash rasters for all five models (all\_to\_all,
         strong heterogeneity, full 0--60\,s).  Each panel shows agent flash
         events as vertical tick marks.  (a) Kuramoto, (b) EAPF consensus,
         (c) EAPF tracker, (d) PCO adaptive PRC, (e) PCO simple.
         Extra flashes from faster agents are visible in panels (d) and (e).}
\label{fig:raster_full}
\end{figure}

% --- Figure 10b: 0-10 s Zoom ---
\begin{figure}[htbp]
\centering
\includegraphics[width=0.95\textwidth]{\figdir/fig10b_representative_flash_rasters_zoom_0_10s.png}
\caption{Representative flash rasters — 0--10\,s zoom showing early
         convergence dynamics.  Kuramoto (a) and EAPF consensus (b) converge
         rapidly; PCO variants (d,e) show extra flashes from the start.}
\label{fig:raster_zoom}
\end{figure}


\FloatBarrier

% ======================================================================
\section{Conclusion}
% ======================================================================

EAPF consensus (event-based phase/frequency consensus with $g_p = 0.02$,
$g_f = 0.02$) is selected as the \textbf{primary HIL candidate} for
Stage 4B under the tested Stage 4A conditions.  Kuramoto ($K = 5.0$) is
selected as the \textbf{secondary comparison model}.  EAPF tracker, PCO
adaptive PRC, and PCO simple are retained as comparative references with
identified limitations.

The primary HIL candidate achieves perfect (100\%) strict zero-lag
synchronisation across all 240 N=3 trials for that model (1200 trials total
across all five models) and all 300 N=5 scalability trials, uses only
flash-event-derived neighbour state (never true internal phases), and has
stable performance across independent random seeds.

\section*{References}

\bibliographystyle{plain}
\bibliography{synchronisation_models/references}

\end{document}
